\section{Optimality and obstruction results}\label{sec:sharpness}

\subsection{A symmetric Laguerre lower-bound family}

We prove sharpness within a symmetric class of densities uniformly comparable with the Gaussian density. The perturbation cancels the first $2m-1$ moments, giving a reflected-MGF defect of factorial size on each fixed interval, while its distribution function differs from the Gaussian one by order $m^{-1/2}$. We use the following Laguerre bound, proved in Appendix~\ref{app:laguerre} from the estimates of Imekraz, Robert and Thomann~\cite[Proposition~3.2]{ImekrazRobertThomann2016}.

\begin{lemma}[Uniform even Laguerre envelope]\label[lemma]{lem:laguerre-envelope}
\begin{equation}\label{eq:CE}
C_E:=\sup_{m\ge2,\,x\in\R}
\sqrt m\,e^{-x^2}\left|L_m^{(-1/2)}(3x^2/2)\right|<\infty.
\end{equation}
\end{lemma}
\begin{proof}
See Appendix~\ref{app:laguerre}.
\end{proof}

Fix $0<\epsilon<1/4$, let $\phi$ be the standard Gaussian density, and define, for $m\ge2$,
\[
h_m(x):=\frac{(-1)^m\sqrt m}{C_E}
 e^{-x^2}L_m^{(-1/2)}(3x^2/2),
\qquad
f_m(x):=\phi(x)(1+\epsilon h_m(x)).
\]
By \eqref{eq:CE}, $|h_m|\le1$, so
\begin{equation}\label{eq:gaussian-comparable-density}
(1-\epsilon)\phi(x)\le f_m(x)\le(1+\epsilon)\phi(x).
\end{equation}
In particular $f_m$ is nonnegative and has exponential moments of every fixed order uniformly in $m$.

The perturbation is even. With $r=3x^2/2$, Laguerre orthogonality for parameter $-1/2$ (Szeg\H{o}~\cite[Chapter~V]{Szego1975}) gives, for $0\le k<m$,
\[
\int_{\R}x^{2k}h_m(x)\phi(x)\,dx
=C_{m,k}\int_0^\infty r^{k-1/2}e^{-r}L_m^{(-1/2)}(r)\,dr=0,
\]
where $C_{m,k}\ne0$ is an irrelevant deterministic factor. The cases $k=0$ and $k=1$ show that $f_m$ has total mass one and variance one, while evenness gives mean zero. Let $W_m$ have density $f_m$.

To compute the MGF exactly, set
\[
J_m(s):=\frac1{\sqrt{2\pi}}
\int_{\R}e^{sx-3x^2/2}L_m^{(-1/2)}(3x^2/2)\,dx.
\]
The Laguerre generating function, first integrated for $|t|$ sufficiently small and then continued analytically to $|t|<1$, gives
\begin{align*}
\sum_{m\ge0}J_m(s)t^m
&=\frac{(1-t)^{-1/2}}{\sqrt{2\pi}}
\int_{\R}
\exp\left(sx-\frac{3x^2}{2(1-t)}\right)dx\\
&=\frac1{\sqrt3}\exp\left(\frac{s^2(1-t)}6\right)\\
&=\frac1{\sqrt3}e^{s^2/6}
\sum_{m\ge0}\frac{(-s^2/6)^m}{m!}t^m.
\end{align*}
Thus
\[
J_m(s)=\frac1{\sqrt3}e^{s^2/6}\frac{(-s^2/6)^m}{m!}.
\]
After multiplying by the factor $(-1)^m\sqrt m/C_E$ in $h_m$ and dividing by the Gaussian MGF $e^{s^2/2}$, we obtain
\begin{equation}\label{eq:even-laguerre-mgf}
M_{W_m}(s)=e^{s^2/2}(1+v_m(s)),
\qquad
v_m(s):=
\frac{\epsilon\sqrt m}{\sqrt3\,C_E}
 e^{-s^2/3}\frac{(s^2/6)^m}{m!}.
\end{equation}
In particular $v_m(s)\ge0$ on the real axis. Since $W_m$ is symmetric,
\begin{equation}\label{eq:even-laguerre-defect}
R_{W_m}(s)=2\log(1+v_m(s))\ge0.
\end{equation}

Fix any $\rho>0$. For $m>\rho^2/3$, the function $s\mapsto e^{-s^2/3}s^{2m}$ is increasing on $[0,\rho]$, since its logarithmic derivative is $2m/s-2s/3>0$ for $0<s\le\rho$. Therefore
\[
\Delta_m:=\Delta_\rho(W_m)
=2\log(1+v_m(\rho)).
\]
Here
\[
v_m(\rho)=\frac{\epsilon\sqrt m}{\sqrt3\,C_E}
 e^{-\rho^2/3}\frac{(\rho^2/6)^m}{m!}\longrightarrow0.
\]
Thus $\Delta_m\asymp v_m(\rho)$, and Stirling's formula yields
\begin{equation}\label{eq:even-laguerre-delta}
\log\frac1{\Delta_m}=m\log m+O_\rho(m).
\end{equation}
Hence
\[
\sqrt{\frac{\log\log(1/\Delta_m)}{\log(1/\Delta_m)}}
\asymp m^{-1/2}.
\]

It remains to prove a Kolmogorov discrepancy of this order. The value at the origin is
\[
L_m^{(-1/2)}(0)=\frac{(1/2)_m}{m!}
=\frac1{4^m}\binom{2m}{m}\asymp m^{-1/2}.
\]
Moreover,
\begin{equation}\label{eq:laguerre-origin-series}
\frac{L_m^{(-1/2)}(r)}{L_m^{(-1/2)}(0)}
=\sum_{k=0}^m
\frac{(-m)_k}{(1/2)_k\,k!}r^k.
\end{equation}
Choose a fixed $c_0>0$ so small that $e^{2c_0}-1\le1/2$. For $0\le r\le c_0/m$, using $|(-m)_k|\le m^k$ and $(1/2)_k\ge2^{-k}k!$ gives
\[
\left|
\frac{L_m^{(-1/2)}(r)}{L_m^{(-1/2)}(0)}-1
\right|
\le\sum_{k\ge1}\frac{(2mr)^k}{(k!)^2}
\le e^{2c_0}-1\le\frac12.
\]
Therefore $L_m^{(-1/2)}(r)$ keeps the sign of its value at the origin and is at least half as large there. Put
\[
x_m:=\sqrt{\frac{2c_0}{3m}}.
\]
For $0\le x\le x_m$, \eqref{eq:CE} and the preceding lower bound show that $h_m(x)$ has the fixed sign $(-1)^m$ and
\[
|h_m(x)|\ge c_1
\]
with $c_1>0$ independent of $m$. Since $f_m-\phi=\epsilon h_m\phi$ is even and has total integral zero, its integral over $(-\infty,0]$ vanishes. Consequently
\[
\begin{aligned}
\dK(W_m,\Normal(0,1))
&\ge |F_{W_m}(x_m)-\Phi(x_m)|\\
&=\epsilon\left|\int_0^{x_m}h_m(x)\phi(x)\,dx\right|
\ge c_2x_m
\ge c_3m^{-1/2}.
\end{aligned}
\]
Together with \eqref{eq:even-laguerre-delta}, this proves \cref{thm:sharpness}. In particular no uniform power-law modulus $C\Delta^\alpha$, $\alpha>0$, can replace the logarithmic modulus in \cref{thm:sharp-upper}.

Since $M_{W_m}(-s)=M_{W_m}(s)>0$ for real $s$, the lower bound persists even when the MGF factor is the positive square root of its reflected product. It is therefore an obstruction to recovery from local transform data, not only an effect of lost Fourier phase.

\subsection{Failure of total-variation stability}

Throughout this subsection, $d_{\mathrm{TV}}(\mu,\nu):=\sup_A|\mu(A)-\nu(A)|$.

\begin{proposition}[Strong-distance obstruction]\label[proposition]{prop:tv}
There is no uniform total-variation stability theorem under the assumptions of \cref{thm:sharp-upper}: a standardized Poisson lattice family has reflected defect tending to zero on every fixed interval while its total-variation distance from every absolutely continuous Gaussian law equals one.
\end{proposition}

\begin{proof}
Let $N_\lambda\sim\operatorname{Poisson}(\lambda)$, with $\lambda\ge1$, and $X_\lambda=(N_\lambda-\lambda)/\sqrt\lambda$. Its centered cumulant generating function is
\[
K_\lambda(s)=\lambda\left(e^{s/\sqrt\lambda}-1-\frac{s}{\sqrt\lambda}\right).
\]
Consequently, uniformly on every fixed compact $s$-interval,
\begin{align*}
R_{X_\lambda}(s)
&=K_\lambda(s)+K_\lambda(-s)-s^2\\
&=2\lambda\left(\cosh(s/\sqrt\lambda)-1-\frac{s^2}{2\lambda}\right)
 =\frac{s^4}{12\lambda}+O_\rho(\lambda^{-2}).
\end{align*}
For fixed $b>0$, the inequality $e^u-1-u\le u^2e^{|u|}/2$ gives
\[
\E e^{b|X_\lambda|}
\le \E e^{bX_\lambda}+\E e^{-bX_\lambda}
\le 2\exp\left(\frac{b^2}{2}e^{b/\sqrt\lambda}\right),
\]
so the family even has a fixed exponential-integrability bound for each fixed $b$. Nevertheless $X_\lambda$ is lattice-valued, whereas $\Normal(0,1)$ is absolutely continuous, so
\[
d_{\mathrm{TV}}(\mathcal L(X_\lambda),\mathcal N(0,1))=1.
\]
\end{proof}

This lattice example separates the metric issue from the sharp Kolmogorov modulus. Related failures of strong-distance stability for Gaussian convolutions, including absolutely continuous symmetric examples, were established by Bobkov, Chistyakov and G\"otze~\cite{BobkovChistyakovGotze2013}. Their regularization results~\cite{BobkovChistyakovGotze2016} show how additional Gaussian smoothing restores total-variation and entropic stability in the convolution setting.
